\documentclass{article}
\usepackage{amsmath}
\usepackage{amssymb}
\usepackage{geometry}
\usepackage{hyperref}
\usepackage{booktabs}
\geometry{a4paper, margin=1in}

\title{Matrix Completion for Causal Inference with Continuous and Reversible Treatments: Theoretical Framework and Implementation Guidelines}
\author{}
\date{}

\begin{document}
\maketitle

\section{Introduction}

The estimation of causal effects in longitudinal and panel data settings has undergone a profound methodological evolution, transitioning away from the rigid structural assumptions of traditional two-way fixed effects (TWFE) and difference-in-differences (DiD) estimators toward highly flexible, imputation-based counterfactual algorithms.[1, 2] Among these modern architectures, matrix completion (MC) methods have emerged as a premier statistical solution for accommodating unobserved, time-varying confounding.[1, 3] By imposing a low-rank structural assumption on the matrix of potential outcomes, these methods elegantly synthesize the horizontal imputation logic of unconfoundedness frameworks with the vertical weighting logic of the synthetic control method (SCM).[1, 4, 5] 

However, the software ecosystems that have democratized these methods---most notably standard R packages such as \texttt{fect}, \texttt{MCPanel}, \texttt{gsynth}, and \texttt{augsynth}---are fundamentally optimized for binary treatment regimes.[6, 7, 8] While recent algorithmic extensions within packages like \texttt{fect} have successfully parameterized reversible (or ``switching'') treatments, where a binary intervention oscillates between active and inactive states over time [2, 6, 9], the explicit integration of a treatment vector that is simultaneously continuous (dosage-based) and reversible remains entirely unsupported in pre-compiled libraries.[6, 10] 

The absence of off-the-shelf programmatic support for continuous, time-varying interventions creates a critical void for researchers analyzing complex dosage dynamics, such as fluctuating carbon tax rates, varying pharmaceutical dosages, dynamic pricing models, or scalable policy implementations.[11, 12, 13] The statistical and computational complexity of this data structure is nontrivial: continuous treatments induce an uncountably infinite number of potential outcomes per unit, of which only a single realization is observed.[14, 15] When compounded by a reversible panel framework---where units can receive varying dosages at different intervals, reduce dosages, or return to a zero-dose baseline---traditional binary unconfoundedness paradigms and standard synthetic control extrapolation mechanisms degrade.[1, 12, 16] 

To overcome the limitations of existing software, analysts must construct bespoke matrix completion estimators from scratch. This report provides an exhaustive, mathematically rigorous, and computationally explicit roadmap for implementing a matrix completion algorithm capable of estimating causal effects for continuous and reversible treatments. By synthesizing foundational low-rank regularization frameworks [1, 17], recent theoretical advancements in factor models for continuous interventions [18, 19, 20], dynamic sequential conditional independence models [21, 22], and state-of-the-art proximal gradient optimization algorithms [23, 24], this analysis establishes the objective functions, numerical solvers, and inference protocols required to programmatically deploy these advanced estimators.

\section{Theoretical Foundations of Causal Matrix Completion}

To understand the algorithmic modifications required for continuous and reversible dosages, it is necessary to first formally define the baseline matrix completion architecture for causal inference, as pioneered by Athey et al. (2021) and extended by the broader econometric literature.[1]

\subsection{The Low-Rank Factor Model for Potential Outcomes}

Consider a longitudinal dataset comprising $N$ units observed over $T$ periods. Let $\mathbf{Y}$ be an $N \times T$ matrix of observed outcomes, where each entry $Y_{it}$ corresponds to the outcome of unit $i$ at time $t$. In a standard binary treatment framework, $W_{it} \in \{0, 1\}$ represents the treatment indicator.[1, 25] The causal inference problem is fundamentally parameterized as a missing data problem: the objective is to accurately impute the missing counterfactual outcomes $Y_{it}(0)$ for the treated units (where $W_{it} = 1$) using the observed control outcomes.[1, 3]

Matrix completion paradigms posit that the potential outcome under the control state, $Y_{it}(0)$, is generated by a multi-factor model heavily influenced by unobserved time-varying confounders:
$$Y_{it}(0) = L_{it} + \alpha_i + \beta_t + X_{it}^\top \gamma + \epsilon_{it}$$
where $\mathbf{L}$ is a low-rank matrix representing interactive fixed effects, $\alpha_i$ and $\beta_t$ are strictly additive unit and time fixed effects, $X_{it}$ is a vector of exogenous time-varying covariates, and $\epsilon_{it}$ is idiosyncratic zero-mean noise.[1, 17, 26] The interactive fixed effects matrix $\mathbf{L}$ is assumed to be the product of unobserved unit-specific factor loadings and time-specific latent factors.[27]

The core structural assumption enabling the imputation of missing values is that the matrix $\mathbf{L}$ is of low rank, meaning it can be decomposed into $\mathbf{U}\mathbf{V}^\top$, where $\mathbf{U}$ is an $N \times R$ matrix and $\mathbf{V}$ is a $T \times R$ matrix, with the rank $R$ being significantly smaller than both $N$ and $T$.[1, 17] This structure is theoretically justified by the economic or biological assumption that a small, finite number of latent, unobserved factors drive the systemic variation across the entire longitudinal panel.[27]

\subsection{Unifying Synthetic Control and Unconfoundedness}

The intellectual leap that established matrix completion as a superior causal estimator is its mathematical unification of the unconfoundedness literature and the synthetic control method.[1, 3] 

Unconfoundedness techniques, such as propensity score matching or unregularized horizontal regression, impute missing values by leveraging similarities across covariates and pretreatment periods for a specific unit, effectively searching for patterns across the rows of the data matrix.[1, 4] Conversely, the synthetic control method operates via vertical regression, identifying a convex combination of untreated donor units (columns) to construct a synthetic counterfactual for the treated entity.[1, 4] 

Matrix completion seamlessly integrates both horizontal and vertical information flows. By regularizing the spectral properties of the outcome matrix through nuclear norm penalization, the algorithm extracts latent commonalities across both the $N$ and $T$ dimensions simultaneously.[2, 12, 26, 28] This bidirectional learning eliminates the need for the strict parallel trends assumption inherent in Difference-in-Differences, as the low-rank matrix organically absorbs non-parallel, unobserved time-varying shocks.[6, 26]

\section{Identifying the Algorithmic Blueprint: Literature Review}

While ready-made R packages do not encompass continuous and reversible matrix completion frameworks [6, 10], an emerging corpus of advanced econometric and computational literature provides the theoretical blueprints necessary to program this functionality from scratch. Building a custom estimator requires synthesizing these disparate mathematical proofs into a single optimization pipeline.

Table 1 delineates the critical scholarly contributions that form the theoretical basis for a scratch implementation.

\begin{table}[h]
\centering
\begin{tabular}{p{0.25\textwidth} p{0.35\textwidth} p{0.35\textwidth}}
\toprule
\textbf{Literature Source} & \textbf{Methodological Innovation} & \textbf{Relevance to Custom Implementation} \\
\midrule
\textbf{Athey et al. (2021)} [1] & Nuclear norm regularization for causal panel data. & Established the fundamental objective function bridging SCM and unconfoundedness via proximal gradient descent. \\
\textbf{Bai \& Ng (2024)} [18, 20, 29] & Counterfactuals in Factor Models with Continuous Treatment. & Formulated the mathematical parameters for average marginal effect (AME) estimation under interactive fixed effects for continuous arrays. \\
\textbf{Poulos (2024)} [12, 30] & Application of MC to continuous regressors (State-Building). & Demonstrated how to embed a continuous treatment variable directly into a nuclear norm regularized least squares objective. \\
\textbf{Xiao \& Wu (2025)} [21, 22] & Continuous treatments with reversibility in dynamic panels. & Formalized the Sequentially Conditional Independence Assumption (SCIA) required for reversible continuous dosage mapping. \\
\textbf{Abadie \& Agarwal (2024)} [27] & Doubly Robust Inference in Latent Factor Models. & Introduced spatial cross-fitting for MC, ensuring finite-sample parametric convergence and debiased imputation. \\
\textbf{Liu, Wang, Xu (2024)} [2, 6] & Architecture for missing data counterfactuals (\texttt{fect}). & Developed essential diagnostic frameworks (carryover tests, placebo equivalence) designed specifically to handle treatment switching. \\
\textbf{Cai et al. (2010)} [24, 31] & Singular Value Thresholding (SVT) Algorithm. & Provided the exact computational algorithm to solve large-scale matrix completion problems without unmanageable matrix inversions. \\
\bottomrule
\end{tabular}
\caption{Critical scholarly contributions}
\end{table}

This literature confirms that while the theoretical boundaries of continuous, reversible causal inference have been solved algebraically, the applied deployment relies heavily on customized algorithmic scripting. To construct this system, the programmer must redefine the causal estimand, rewrite the objective function, select an appropriate iterative solver, and design a custom inference protocol.

\section{The Challenge of Continuous and Reversible Treatments}

Transitioning from the foundational binary model to one supporting continuous and reversible treatments requires addressing distinct econometric challenges. The standard binary masking matrix $\mathbf{W}$, utilized in conventional algorithms, is structurally insufficient when the treatment indicator $D_{it} \in \mathbb{R}$.[6, 8] 

\subsection{Defining the Continuous, Reversible Estimand}

Let $D_{it}$ denote a continuous treatment dosage applied to unit $i$ at time $t$. The treatment is reversible, meaning $D_{it}$ can fluctuate continuously over time, including returning to $0$ (the defined baseline or pure control state).[6, 21, 32] The potential outcome function is mapped continuously rather than discretely: $Y_{it}(d)$ represents the theoretical outcome for unit $i$ at time $t$ if the unit were exposed to dosage level $d$.[14, 33]

The causal estimands in this paradigm necessarily shift away from the simple Average Treatment Effect on the Treated (ATT). Because ``treatment'' is a continuum, researchers must target the Average Marginal Effect (AME), the Average Causal Response (ACR), or the full continuous dose-response function.[5, 11, 20, 34] 

If the treatment effect is modeled parametrically and additively, the observed outcome is defined as:
$$Y_{it} = Y_{it}(0) + \tau_{it}(D_{it})$$
where $\tau_{it}(\cdot)$ represents the heterogeneous dose-response function corresponding to the dosage $D_{it}$.[35, 36] If the effect is assumed strictly linear, this simplifies to $\phi \cdot D_{it}$, where $\phi$ is the constant marginal effect. However, modern implementations favor non-parametric extraction of the dose-response to account for diminishing returns or threshold effects.[37, 38]

\subsection{Sequentially Conditional Independence and Dynamic Confounding}

When a continuous treatment reverses over time, researchers must rigorously account for sequential endogeneity. In dynamic real-world environments, current outcomes frequently depend on historical treatments, and future continuous treatments may be iteratively assigned based on past outcomes.[22] 

Xiao and Wu (2025) demonstrate that identifying valid causal effects with a continuous, time-varying treatment requires invoking a Sequentially Conditional Independence Assumption (SCIA).[21, 22] SCIA ensures that, conditional on the history of treatments and observed covariates, the unobserved potential outcomes are independent of the current continuous treatment assignment.[21, 22] 

Furthermore, reversible continuous treatments introduce the severe risk of ``carryover effects,'' where the biological or economic impact of a previous high dosage $D_{i, t-1}$ persists into period $t$ even if the current dosage $D_{it}$ is reduced or returned to $0$.[6, 32] The assumption of strict exogeneity is frequently invoked in basic counterfactual estimators to prohibit these feedback loops. To implement a robust matrix completion algorithm for this specific data structure from scratch, the computational architecture must systematically separate the estimation of the baseline potential outcome matrix $\mathbf{Y}(0)$ from the estimation of the structural dose-response function $\tau(D_{it})$ to prevent lagged dosage effects from corrupting the latent factor imputation.

\section{Mathematical Formulation of the Estimator}

To implement the algorithm programmatically, the estimation problem must be formulated as a convex optimization problem. The primary objective is to accurately isolate the low-rank component $\mathbf{L}$ using only those observations where the treatment effect is strictly zero (i.e., $D_{it} = 0$), or by jointly estimating the low-rank matrix alongside the continuous dose-response parameters over the entire observation matrix.[5, 39]

\subsection{The Nuclear Norm Regularized Objective Function}

Let $\Omega$ represent the set of indices $(i,t)$ where the unit is in the pure control state ($D_{it} = 0$). For these specific observations, the observed outcome is simply the baseline counterfactual $Y_{it} = Y_{it}(0)$. For units receiving any continuous dosage ($D_{it} > 0$ or $D_{it} < 0$), the baseline outcome is treated algorithmically as missing or systematically adjusted.[25] 

If assuming a global parametric dose-response curve, the potential outcome under treatment is defined as $Y_{it} = L_{it} + \alpha_i + \beta_t + X_{it}^\top \gamma + \phi(D_{it}, \theta) + \epsilon_{it}$, where $\phi(D_{it}, \theta)$ is the functional effect of the continuous treatment parametrized by the vector $\theta$.[12, 40]

The optimization problem seeks to minimize the regularized least squares objective over the entire matrix, utilizing the Nuclear Norm $\| \mathbf{L} \|_*$ to penalize rank complexity and prevent overfitting.[12, 41] The nuclear norm, defined algebraically as the sum of the singular values of the matrix ($\|\mathbf{L}\|_* = \sum_k \sigma_k(\mathbf{L})$), serves as the tightest convex relaxation of the non-convex rank minimization problem, ensuring that a global minimum can be found via standard descent methods.[41, 42, 43]

The full joint objective function to be programmed is:
$$\min_{\mathbf{L}, \alpha, \beta, \gamma, \theta} \mathcal{J} = \frac{1}{|\Omega_{all}|} \sum_{i=1}^N \sum_{t=1}^T \left( Y_{it} - L_{it} - \alpha_i - \beta_t - X_{it}^\top \gamma - \phi(D_{it}, \theta) \right)^2 + \lambda_L \|\mathbf{L}\|_*$$

\subsection{Decoupling the Objective for Reversible Treatments}

In settings with complex reversible treatments where assuming a global parametric dose-response function $\phi$ is too restrictive, the algorithm should be mathematically decoupled into a two-stage imputation pipeline. This mirrors the internal logic of the \texttt{fect} package's counterfactual generation but applies it to the continuous continuum.[6]

\textbf{Stage 1: Imputation of the Baseline Matrix via Masking}
Construct a binary observation mask matrix $\mathbf{M}$, where $M_{it} = 1$ if $D_{it} = 0$ (the unexposed, baseline control state), and $M_{it} = 0$ if $D_{it} \neq 0$ (active continuous treatment state).[25] 

The custom algorithm solves the restricted nuclear norm minimization problem exclusively over the $\Omega$ subset (where $M_{it} = 1$):
$$\min_{\mathbf{L}, \alpha, \beta, \gamma} \frac{1}{|\Omega|} \sum_{(i,t) \in \Omega} \left( Y_{it} - L_{it} - \alpha_i - \beta_t - X_{it}^\top \gamma \right)^2 + \lambda_L \|\mathbf{L}\|_*$$

This step requires careful computational handling of the reversibility. Because the continuous dosage $D_{it}$ switches on and off repeatedly, the observed baseline entries $\Omega$ will not be restricted strictly to neat pre-treatment blocks; they will be scattered sporadically throughout the matrix.[6, 44] 

Paradoxically, this sporadic distribution improves the theoretical convergence guarantees of the matrix completion algorithm. The fundamental theorem of exact matrix completion (established by Candès \& Recht, 2009) relies heavily on the ``incoherence'' assumption---the requirement that missing entries are somewhat uniformly or randomly distributed.[1, 45, 46] Reversible treatments naturally approximate this mathematical condition far better than traditional block-staggered designs.[47] The matrix completion algorithm leverages observed post-treatment reversions to zero to heavily calibrate the unit-specific factor loadings, making it vastly superior to synthetic control methods which rely almost exclusively on continuous pre-treatment blocks.[4, 37, 47, 48]

\textbf{Stage 2: Calculation of the Continuous Residual}
Once the full low-rank matrix $\mathbf{L}$ is imputed via optimization, the estimated baseline counterfactual $\hat{Y}_{it}(0)$ is constructed for all treated units:
$$\hat{Y}_{it}(0) = \hat{L}_{it} + \hat{\alpha}_i + \hat{\beta}_t + X_{it}^\top \hat{\gamma}$$

The individual, time-specific causal effect is cleanly isolated as the residual difference:
$$\hat{\tau}_{it} = Y_{it} - \hat{Y}_{it}(0)$$

With a continuous treatment, these heterogeneous treatment effects $\hat{\tau}_{it}$ are subsequently analyzed against the continuous dosage $D_{it}$ to estimate the dose-response curve.

\section{Algorithmic Design: Optimization via Proximal Gradient Solvers}

To implement the objective function from scratch without relying on compiled packages like \texttt{MCPanel}, the programmer must select an appropriate numerical solver. Standard Ordinary Least Squares (OLS) or Alternating Least Squares (ALS) fail when penalizing the nuclear norm due to the non-differentiability of the singular values at zero.[49, 50] 

The most stable and computationally efficient numerical solver for this specific architecture is the Iterative Soft Thresholding Algorithm (ISTA), specifically its matrix variant known as Singular Value Thresholding (SVT) or the Soft-Impute algorithm.[23, 24, 31, 37, 51] More advanced parallelized implementations may utilize the Preconditioned Iterative Soft Thresholding Algorithm (pISTA) or the Alternating Direction Method of Multipliers (ADMM).[23, 43, 50]

\subsection{The Proximal Gradient Method and SVT}

The objective function consists of a smooth, continuously differentiable loss component (the squared error over the observed data) and a non-smooth, convex penalty (the nuclear norm).[50, 52, 53] Proximal gradient descent is the standard numerical technique designed explicitly for such composite functions.

Let $f(\mathbf{L})$ represent the least squares loss over the observed entries $\Omega$. The gradient of the loss with respect to $\mathbf{L}$ is evaluated as:
$$\nabla f(\mathbf{L})_{it} = \begin{cases} 2 (L_{it} - \tilde{Y}_{it}) & \text{if } (i,t) \in \Omega \\ 0 & \text{if } (i,t) \notin \Omega \end{cases}$$
where $\tilde{Y}_{it} = Y_{it} - \alpha_i - \beta_t - X_{it}^\top \gamma$ is the residualized outcome.[17, 41]

The proximal operator for the nuclear norm penalty $\lambda_L \|\mathbf{L}\|_*$ is the Singular Value Soft Thresholding operator $\mathcal{S}_{\lambda_L}(\cdot)$.[24] For any given matrix $\mathbf{A}$ with a Singular Value Decomposition (SVD) defined as $\mathbf{A} = \mathbf{U} \mathbf{\Sigma} \mathbf{V}^\top$, the thresholding operator shrinks the singular values toward zero:
$$\mathcal{S}_{\lambda_L}(\mathbf{A}) = \mathbf{U} \mathbf{\Sigma}_{\lambda_L} \mathbf{V}^\top$$
where the diagonal matrix $\mathbf{\Sigma}_{\lambda_L}$ contains the thresholded entries $\max(\sigma_k - \lambda_L, 0)$.[31, 54, 55] Any singular value smaller than the tuning parameter $\lambda_L$ is forced to exactly zero, guaranteeing that the resulting matrix strictly adheres to a low-rank structure.[41, 56]

\subsection{The Soft-Impute Algorithm (Implementation Blueprint)}

To implement this from scratch in R, Python, or Julia, the Soft-Impute iterative procedure operates via the following programmed steps [57]:

\begin{enumerate}
    \item \textbf{Initialization Phase}: 
    \begin{itemize}
        \item Center the observed outcome matrix $\mathbf{Y}$ by estimating and removing simple unit ($\alpha_i$) and time ($\beta_t$) fixed effects from the $\Omega$ subset.[17]
        \item Initialize the covariate parameter vector $\gamma^{(0)}$ to zero.
        \item Initialize the low-rank matrix $\mathbf{L}^{(0)} = \mathbf{0}$.
        \item Define a sequence of penalty values $\lambda_L$ for cross-validation, typically starting from the largest singular value of the zero-filled matrix and descending logarithmically.[58]
    \end{itemize}

    \item \textbf{Iterative Update Loop} (Executed for a fixed $\lambda_L$):
    \begin{itemize}
        \item \textbf{Step A (Gradient Construction / Matrix Completion)}: Construct the complete dense matrix $\mathbf{Z}^{(k)}$ by replacing the missing entries (where continuous treatment is active) in the observed matrix with the current estimates from the previous iteration $\mathbf{L}^{(k-1)}$:
        $$Z_{it}^{(k)} = \begin{cases} Y_{it} - X_{it}^\top \gamma^{(k-1)} & \text{if } D_{it} = 0 \\ L_{it}^{(k-1)} & \text{if } D_{it} \neq 0 \end{cases}$$
        \item \textbf{Step B (Singular Value Decomposition)}: Perform the SVD on the newly completed matrix $\mathbf{Z}^{(k)} = \mathbf{U} \mathbf{\Sigma} \mathbf{V}^\top$. In R, implementing this via truncated SVD packages (such as \texttt{irlba}) rather than base \texttt{svd()} is critical for computational feasibility on large panel data.[42]
        \item \textbf{Step C (Soft Thresholding Application)}: Apply the soft-thresholding mathematical operator to the singular values:
        $$\mathbf{L}^{(k)} = \mathbf{U} \max(\mathbf{\Sigma} - \lambda_L, 0) \mathbf{V}^\top$$
        \item \textbf{Step D (Covariate Coefficient Update)}: Conditional on the newly updated $\mathbf{L}^{(k)}$, update the continuous covariate coefficients $\gamma^{(k)}$ using standard Ordinary Least Squares executed strictly on the untreated subset $\Omega$:
        $$\gamma^{(k)} = \left( \sum_{(i,t) \in \Omega} X_{it} X_{it}^\top \right)^{-1} \sum_{(i,t) \in \Omega} X_{it} (Y_{it} - L_{it}^{(k)})$$
    \end{itemize}

    \item \textbf{Convergence Verification}: At the end of each loop, evaluate the Frobenius norm of the difference between successive iterations: 
    $$\frac{\|\mathbf{L}^{(k)} - \mathbf{L}^{(k-1)}\|_F^2}{\|\mathbf{L}^{(k-1)}\|_F^2}$$
    If this calculated ratio falls below a specified precision tolerance threshold (e.g., $10^{-5}$), terminate the iterative loop.[52]
\end{enumerate}

Table 2 highlights the specific programmatic algorithms and solver variations that can be implemented depending on the computational environment.

\begin{table}[h]
\centering
\begin{tabular}{p{0.3\textwidth} p{0.35\textwidth} p{0.3\textwidth}}
\toprule
\textbf{Algorithm Variant} & \textbf{Optimization Mechanics} & \textbf{Ideal Use Case for Implementation} \\
\midrule
\textbf{Singular Value Thresholding (SVT)} [24, 31] & Alternates dense matrix completion with exact SVD soft-thresholding. & Standard panel datasets ($N, T < 5000$). Easiest to code in base R/Python. \\
\textbf{Alternating Direction Method of Multipliers (ADMM)} [43] & Uses augmented Lagrangians to decouple the loss and the nuclear norm penalty. & Complex constrained models where covariates $X_{it}$ are high-dimensional. \\
\textbf{Preconditioned ISTA (pISTA)} [23, 50] & Uses the inverse Hessian as a preconditioner to accelerate convergence rates. & Massive panel arrays requiring GPU-accelerated parallelization architectures. \\
\bottomrule
\end{tabular}
\caption{Optimization algorithm variants}
\end{table}

This ISTA/Soft-Impute framework deliberately bypasses the complex Hessian inversions required by second-order Newton algorithms, offering highly parallelizable operations uniquely suitable for large-scale economic and biological panel data.[23, 50] 

\section{Estimation of Causal Estimands and Dose-Response Functions}

With the baseline counterfactual matrix $\hat{\mathbf{Y}}(0)$ successfully imputed via the custom SVT routine, the subsequent analytical stage involves aggregating the unit-time causal effects into meaningful continuous estimands and extracting the structural dose-response function.[5, 11]

\subsection{Average Marginal Effects and Average Causal Response}

Because the treatment is continuous, traditional binary estimators like the ATT are mathematically ill-defined without arbitrarily specifying a binary dosage boundary. Instead, the relevant estimands are the Average Causal Response (ACR) or the Average Marginal Effect (AME).[20, 21, 22]

If modeling the structural dose-response linearly, the researcher can compute the Average Marginal Effect by projecting the isolated heterogeneous treatment residuals $\hat{\tau}_{it} = Y_{it} - \hat{Y}_{it}(0)$ onto the continuous dosage matrix $D_{it}$ [20, 29]:
$$\hat{\phi}_{AME} = \frac{\sum_{i,t} W_{it} D_{it} \hat{\tau}_{it}}{\sum_{i,t} W_{it} D_{it}^2}$$
where $W_{it}$ is an indicator that the dosage $D_{it} > 0$. 

However, assuming linearity across a continuous treatment continuum is frequently overly restrictive. A superior non-parametric approach involves mapping the dose-response curve using flexible machine learning smoothers. The individual treatment effects $\hat{\tau}_{it}$ are utilized as the dependent variable, and restricted cubic splines, generalized additive models (GAMs), or generalized random forests are regressed on $D_{it}$ to non-parametrically uncover threshold effects and diminishing returns.[34] 

\subsection{Doubly Robust Weighting for Continuous Treatments}

To further immunize the custom matrix completion estimator against model misspecification, the algorithm can be enhanced by integrating Inverse Probability Weighting (IPW), creating a doubly robust estimator.[27] In purely observational settings with continuous treatments, certain high dosages may be heavily correlated with specific unobserved latent factors or observed covariates.

The standard proximal gradient objective function can be mathematically augmented by incorporating a weighting matrix $\mathbf{\Omega}^{IPW}$. This matrix is derived from a Generalized Propensity Score (GPS) model explicitly tailored for continuous treatments.[15, 59] The GPS assesses the conditional probability density of observing the specific continuous treatment level given the covariates, $f(D_{it} | X_{it})$.

Recent advancements demonstrate that using the Covariate Balancing Generalized Propensity Score (CBGPS) for continuous scenarios dramatically improves finite-sample performance.[60] CBGPS optimizes the propensity scores by directly minimizing the correlation between the continuous treatment and the confounding covariates, rather than relying on maximum likelihood estimation.[60]

The modified, doubly robust objective function becomes:
$$\min_{\mathbf{L}, \beta} \frac{1}{|\Omega_{all}|} \sum_{i,t} \omega_{it} \left( Y_{it} - L_{it} - X_{it}^\top \beta - \phi(D_{it}) \right)^2 + \lambda_L \|\mathbf{L}\|_*$$
where $\omega_{it}$ represents the stabilized continuous propensity scores or entropy balance weights.[14, 15, 26] This doubly robust formulation guarantees that the causal estimator remains consistent and unbiased even if the low-rank assumption governing $\mathbf{L}$ is mildly misspecified, provided that the continuous treatment assignment mechanism is modeled correctly via the weights.[27]

\section{Hyperparameter Tuning, Inference, and Diagnostic Protocols}

The statistical integrity of any regularized causal model relies heavily on appropriate hyperparameter tuning, rigorous quantification of uncertainty, and aggressive diagnostic testing of structural assumptions. Programming these validation protocols from scratch requires specialized adaptations for time-series cross-sectional data.

\subsection{Cross-Validation for Serial Correlation}

The predictive accuracy of the matrix completion estimator hinges almost entirely on the selection of the regularization parameter $\lambda_L$, which dictates the rank and complexity of the imputed latent matrix.[1, 17, 61] To program this tuning phase from scratch, one must deploy a structured cross-validation scheme.

Given the time-series cross-sectional nature of the panel data, standard random hold-out $K$-fold validation is invalid, as it violates inherent temporal dependence structures and auto-correlation. Instead, block cross-validation or leave-one-period-out cross-validation must be utilized.[58] 

The custom script should execute the following sequence:
\begin{enumerate}
    \item Partition the observed control set $\Omega$ into $K$ spatial or temporal folds, blocking strictly by time periods or by intact unit trajectories to respect serial correlation.
    \item For each $\lambda_L$ in a predefined logarithmic grid, execute the Soft-Impute algorithm on the remaining $K-1$ folds.
    \item Calculate the predictive Mean Squared Error (MSE) strictly on the hold-out fold.
    \item Select the specific $\lambda_L$ that globally minimizes the cross-validated MSE to prevent both over-shrinkage (bias) and under-shrinkage (variance).
\end{enumerate}

\subsection{Non-Parametric Block Bootstrapping for Inference}

Traditional asymptotic inference for matrix completion is notoriously complex. Because the nuclear norm is non-differentiable at the origin, standard derivations of asymptotic normality fail.[50] While highly advanced theoretical work by Abadie and Agarwal (2024) introduces cross-fitted spatial split-sample procedures to establish parametric convergence rates mathematically [27], implementing this architecture from scratch is often prohibitively dense for applied research. 

A far more tractable programmatic approach, consistently utilized in the continuous treatment econometric literature [22, 62], is the non-parametric block bootstrap:
\begin{enumerate}
    \item Programmatically resample units $i \in \{1, \dots, N\}$ with replacement to generate a bootstrap sample dataset. Crucially, the algorithm must preserve the entire time series $T$ intact for each sampled unit to accurately reflect within-unit serial correlation and heteroskedasticity.[22]
    \item Re-run the full cross-validation grid search, the Soft-Impute convergence algorithm, and the continuous dose-response extraction on the newly created bootstrap sample.
    \item Repeat this exhaustive process $B$ times (e.g., $B=500$) to successfully generate a robust empirical distribution of the causal estimand (e.g., the AME).
    \item Derive 95\% confidence intervals directly from the 2.5th and 97.5th percentiles of the generated bootstrap distribution.
\end{enumerate}

\subsection{Diagnostic Testing for Continuous Reversibility}

When treatments are continuously reversible, the strict exogeneity assumption must be aggressively tested to ensure that unmodeled feedback loops are not biasing the low-rank imputation. The diagnostic architecture developed for the \texttt{fect} package establishes a critical precedent for equivalence testing that should be directly programmed into the custom matrix completion algorithm.[2, 6]

Table 3 details the structural diagnostic tests required for validation.

\begin{table}[h]
\centering
\begin{tabular}{p{0.25\textwidth} p{0.35\textwidth} p{0.35\textwidth}}
\toprule
\textbf{Diagnostic Protocol} & \textbf{Theoretical Purpose} & \textbf{Algorithmic Implementation Strategy} \\
\midrule
\textbf{Carryover Effect Test} [6, 32] & Verifies that a continuous treatment does not have lingering unmodeled effects after the dosage returns to zero. & Identify ``wash-out'' periods where $D_{i, t-1} > 0$ but $D_{it} = 0$. Algorithmically mask these periods during imputation. Evaluate if $\hat{Y}_{it}(0)$ significantly deviates from observed $Y_{it}$. \\
\textbf{Pre-trend Placebo Test} [6] & Confirms that the latent factor model effectively captures unobserved confounding prior to dosage exposure. & Artificially mask a set of pure control periods prior to the first non-zero dosage. Impute the counterfactuals and test if the residual $\hat{\tau}_{it}^{placebo}$ is statistically indistinguishable from zero. \\
\textbf{DAG Verification} [21, 22] & Tests the validity of the Sequentially Conditional Independence Assumption (SCIA). & Translate contextual knowledge into conditional independence relationships through formal algorithmic deduction and covariate lagged screening. \\
\bottomrule
\end{tabular}
\caption{Structural diagnostic tests}
\end{table}

If the carryover test fails---indicating that the estimated counterfactual $\hat{Y}_{it}(0)$ diverges from the observed reality precisely when the continuous treatment is paused---the model suffers from severe carryover bias. The custom algorithm must then be structurally updated to include lagged dosages $D_{i, t-1}, D_{i, t-2}$ directly into the fixed covariate matrix $X_{it}$ to absorb the persistent effect.[6]

\section{Conclusion}

The ongoing evolution of causal inference in longitudinal panel data requires methodologies that can gracefully handle complex, real-world intervention dynamics, specifically encompassing continuous and reversible dosage treatments. Current programmatic libraries in R and Python, while highly optimized for binary and staggered designs, lack the requisite algebraic flexibility to manage these complex data structures natively. By constructing a custom matrix completion estimator from scratch, researchers can seamlessly accommodate time-varying dosages and switching mechanisms without sacrificing theoretical rigor.

This scratch implementation requires a fundamental transition in the estimation paradigm. The reliance on discrete treatment masks must be replaced with continuous dose-response mapping, heavily relying on the Proximal Gradient and Singular Value Thresholding (SVT) algorithms to solve the nuclear norm regularized objective function efficiently. By methodically isolating the estimation of the baseline, low-rank unconfounded matrix from the evaluation of the continuous treatment residual, this custom methodology effectively circumvents the structural limitations of standard software packages. When meticulously paired with block-bootstrapped standard errors, continuous Covariate Balancing Generalized Propensity Score weighting, rigorous carryover diagnostics, and dynamic sequential conditional independence modeling, the bespoke algorithmic implementation detailed in this analysis provides a mathematically rigorous, scalable framework for evaluating continuous, reversible causal effects in highly confounded, high-dimensional environments.

\end{document}